\section{Summary} \label{sec:summary}
% ============================================================================
In this paper, we performed time-resolved empirical spectroscopy of $106$
single-pulse, bright Fermi/GBM bursts ($1057$ spectra; six models each), one
of the largest uniform single-pulse spectral surveys to date. Our main findings
are as follows.
\begin{enumerate}
\item The low-energy index peaks at $\alpha\approx-0.7$, near the slow-cooled
  synchrotron limit, with $74\%$ ($74/100$) of bursts reaching harder values in
  at least one bin; the rise-phase \Ep\ evolution is HTS-dominated
  ($63\%$, $48/76$), as in previous single-pulse samples.
\item Of the bins requiring curvature beyond a single break, $91\%$ ($151/166$)
  are equally or better described by adding a blackbody than by an explicit
  second break (the $9\%$ two-break fraction is a lower limit); no single model
  is decisively preferred over the next-best in $92\%$ ($875/950$) of
  comparable bins.
\item The \citet{Burgess2014a} $\Ep$--$\kT$ correlation is recovered in $89\%$
  ($17/19$) of the bursts with enough decisive-blackbody bins to test
  (GRB~130427A: $\rho=0.93$, $\Ep\propto\kT^{0.88\pm0.07}$), but fainter bursts
  cannot be tested with empirical models---a model-sensitivity, not a sample,
  limit.
\item The two 2SBPL break energies are positively correlated ($\rho=0.57$;
  \citealt{DAgostini2005} slope $0.46\pm0.05$, \ssc$=0.38$~dex), both within
  and between bursts.
\end{enumerate}
We conclude that the dominant spectral curvature of bright single-pulse GRBs is
more economically described by a sub-dominant photospheric component than by a
second synchrotron break, with a genuine two-break minority---epitomized by
GRB~110721A---in which the cooling and injection frequencies are correlated.

